
\chapter{Balance Laws}

%\newcommand{\phibar}{\bar{\phi}}
%\newcommand{\phistar}{\overset{*}{\phi}}
\newcommand{\phistar}{\phi^*}
\newcommand{\phihat}{\hat{\phi}}

In this section, we posit the {\bf Master Balance Law} 
in general form.  Also called ``conservation of anything,'' this
description assumes an arbitrary control volume $V \in \realthree$
with boundary $\partial V$.  What follows is the so-called 
Eulerian formulation, illustrated in Fig.~\ref{fig:mass_balance}, 
which assumes the control volume, 
while arbitrary in its placement in space, is {\bf fixed} 
in space.\footnote{The connection to the Lagrangian formulation, 
in which the control is not fixed in space, but instead, 
moves through space as the body to which it
is attached moves through space, is examined in the sequel.} 
\begin{figure}[htb]
  \begin{center}
    %\begin{overpic}[width=0.85\textwidth, grid, tics=10]{fig/mass_balance}
    \begin{overpic}[width=0.85\textwidth,natwidth=566,natheight=283]{fig/mass_balance.pdf}
      \put(22, 36){flux in}
      \put(85, 26){flux out}
      \put(44,  2){$V$}
      \put(54.5, 47){$\vnhat$ on $\partial V$}
      %\put(50.5,  24){$g$}
      \put(49.5,  23.8){$\phi^*$}
      %
    \end{overpic}
  \end{center}
  \caption{Concept of a fixed control volume of the Eulerian formulation.}
  \label{fig:mass_balance} 
\end{figure}

Now, let 
\begin{itemize}
\item $\phi(\vx, t)$ be the description of some physical quantity 
per unit volume, which can in general be a scalar, vector, or tensor quantity, 

\item $\phistar(\vx, t)$ be the generation (depletion) {\em rate} of the 
physical quantity per unit volume, 

\item $\phihat(\vx, t)$ be the physical quantity per unit volume 
that flows across the boundary with rate (velocity) $\vv(\vx, t)$, and 

\item $\partial V$ be the boundary of the control volume, 
with outward pointing surface normal unit vector $\vnhat(\vx, t)$.  
\end{itemize} 
The time rate of change of the quantity $\phi(\vx,t)$ within the 
differential control volume $dV$ is then equal to 
\begin{itemize}
\item {\bf generation}---the rate at which the quantity is generated (or depleted) 
within the control volume by the point source (point sink), plus
\item {\bf flux}---the rate at which the quantity is either carried into (or out of) 
the volume across the volume's differential surface area $dA$.
\end{itemize}
These words, stated mathematically, are
\be
 \underset{\mbox{\footnotesize rate}}{\underbrace{\frac{d}{dt} 
 \int_{V} \phi(\vx,t) \; dV}}
  \;\; = \;\;  
   \underset{\mbox{\footnotesize source}}{\underbrace{\int_{V} 
   \phistar(\vx,t) \; dV}}
   \;\; - \;\;
   \underset{\mbox{\footnotesize flux}}{\underbrace{\int_{\partial V} 
   %\left(
   \phihat(\vx,t) \; \vv(\vx,t) 
   %\right) 
   \cdot \vnhat \; dA}}
   ,
\ee
which is the general form of the 
{\bf Master Balance Law}.  We adopt the following sign convention 
to impart consistency:
\begin{itemize}
\item For positive time rate of change of the quantity $\phi(\vx,t)$, 
the generation $\phistar(\vx,t)$ is trivially positive.  
For quantity $\phihat(\vx,t)$ carried
{\em into} the control volume, the inner product of $\vv$ with $\vnhat$ is
negative, thus a negative outside of the flux integral is used.  
\item Likewise, for a negative time rate of change of 
quantity $\phi(\vx,t)$, 
the generation $\phistar(\vx,t)$ is negative.  
And, for the quantity $\phihat(\vx,t)$ 
carried {\em out of} the control volume, the inner product of $\vv$ with
$\vnhat$ is positive, thus the negative outside of the flux integral 
provides sign consistency.  
\end{itemize}

\section{Balance Law Metaphor}
Here we present the Master Balance Law by way of metaphor.  
Let the volume $V$ be the geographic region of Oregon, 
with boundary $\partial V$ that touches the state borders of
Washington, Idaho, Nevada, and California 
in the United States.\footnote{In this oversimplified
metaphor, we assume flux does not occur via the western seaboard nor via the air.} 
Let $\phi$ be the number of people living in the state of Oregon.
The change in population $\Delta \phi$ over some time interval $\Delta t$ 
is equal to the 
net generation rate $\phistar$ (number of births less the number of deaths
in the time interval), plus the net population flux rate $\phihat \vv$
(the number of people moving to Oregon less
the number of people moving from Oregon within the time interval).

\section{Conservation of Mass}
Let us now apply the Master Balance Law to the physical 
quantity of mass, and develop a relationship for 
{\bf conservation of mass},
a concept where mass is neither created nor destroyed.  
The physical quantity of mass per unit volume is 
density $\rho(\vx,t)$.  Thus, we substitute $\rho(\vx,t)$ 
for $\phi(\vx,t)$ in the foregoing Master Balance Law.

\be
\begin{tabular}{cl|l}
 &
 {\bf Assumption}
 &
 {\bf Expression} \\ \hline \hline
 (1) 
 &
 No mass generation or depletion. 
 &
 $\rho^*(\vx,t) = 0$
 \\ \hline
 (2)
 &
 Mass can flow in and out of the boundaries.
 &
 $\rho(\vx, t) \vv(\vx, t) \cdot \hat{\vn}(\vx, t) \neq 0$
\end{tabular}
\ee

\begin{Hresult}
Conservation of mass in the current configuration is stated as 
\be
   \frac{\partial \rho(\vx,t)}{\partial t} + \grad \cdot 
   %\left[ \; \rho(\vx,t) \;  \vv(\vx,t) \; \right] 
   \left( \; \rho(\vx,t) \;  \vv(\vx,t) \; \right)
 = 
 0
 \;\; \Longleftrightarrow \;\; \rho_{,t} + (\rho \; v_{i})_{,i} = 0.
 \label{eq:mass_conservation}
\ee
\end{Hresult}

\begin{Hproof}
For mass to be conserved, 
mass within a control volume is neither created nor destroyed, thus the source
term is identically zero.  Thus, increases (decreases) in the mass rate 
term are produced only by flux of mass in through (out from) the boundary. 
The idea, stated mathematically, is
\be
 \underset{\mbox{\footnotesize rate}}{\underbrace{\frac{d}{dt} 
 \int_{V} \rho(\vx,t) \; dV}}
  \;\; = \;\;  
   \underset{\mbox{\footnotesize source = 0}}{\underbrace{\int_{V} 
   \rho^*(\vx,t) \; dV}}
   \;\; - \;\;
   \underset{\mbox{\footnotesize flux}}{\underbrace{\int_{\partial V} 
   %\left[\; \rho(\vx,t) \; \vv(\vx,t) \; \right] \cdot \vnhat(\vx,t) \; dA}}
   \rho(\vx,t) \; \vv(\vx,t) \cdot \vnhat(\vx,t) \; dA}}
   .
\ee
The Divergence Theorem is used to write the flux term surface integral 
as a volume integral, thus
\be
 \underset{\mbox{\footnotesize rate}}{\underbrace{\frac{d}{dt} 
 \int_{V} \rho(\vx,t) \; dV}}
  \;\; = \;\;  
   0
   \;\; - \;\;
   \underset{\mbox{\footnotesize flux}}{\underbrace{\int_{V} 
   %\grad \cdot \left[\;\rho(\vx,t) \; \vv(\vx,t) \; \right] \; dV}}
   \grad \cdot \left(\;\rho(\vx,t) \; \vv(\vx,t) \; \right) \; dV}}
   .
\ee
Because $V$ is fixed in an inertial frame, the foregoing equation can be 
rewritten with the rate time derivative inside the integral, thus
\be
 \int_{V} \left(
   \frac{\partial \rho(\vx,t)}{\partial t} + \grad \cdot 
   %\left[ \; \rho(\vx,t) \;  \vv(\vx,t) \; \right] \; \right) \; dV 
   \left( \; \rho(\vx,t) \;  \vv(\vx,t) \; \right) \; \right) \; dV 
 = 
 0
 .
\ee
Finally, arbitrariness of $B_t$ 
requires that the integrand is identically zero, thus
\be
   \frac{\partial \rho(\vx,t)}{\partial t} + \grad \cdot 
   %\left[ \; \rho(\vx,t) \;  \vv(\vx,t) \; \right] 
   \left( \; \rho(\vx,t) \;  \vv(\vx,t) \; \right)
 = 
 0
 \;\; \Longleftrightarrow \;\; \rho_{,t} + (\rho \; v_{i})_{,i} = 0,
 \label{eq:mass_conservation_two}
\ee
which is the result.
\end{Hproof}


\begin{Hremark}
Conservation of mass is also called the {\bf continuity equation}.  
Some texts will recast the continuity equation by defining a 
{\bf material derivative} of density as follows:  Since
$\rho = \rho(\vx, t)$, then 
\be
 \dot \rho = \frac{\partial \rho}{\partial \vx} \frac{\partial \vx}{\partial t} 
 + \frac{\partial \rho}{\partial t} \defe \frac{D \rho}{Dt}.
\ee
With arguments $(\vx, t)$ suppressed for brevity, (\ref{eq:mass_conservation}) 
can be expanded to give 
\be
  \frac{\partial \rho}{\partial t} + 
  \grad \rho \cdot \vv + \rho \; \grad \cdot \vv = 
  \underset{\defe \; \dot\rho}{\underbrace{
  \frac{\partial \rho}{\partial t} + 
  \frac{\partial \rho}{\partial \vx} \frac{\partial \vx}{\partial t}}} +
  \rho \grad \cdot \vv = 0,
\ee
or simply
\be
\dot \rho + \rho \grad \cdot \vv = 0.
\ee
\end{Hremark}


\section{Conservation of Linear Momentum}

\section{Conservation of Angular Momentum}

\be
\begin{tabular}{cl|l}
 &
 {\bf Assumption}
 &
 {\bf Expression} \\ \hline \hline
 (1) 
 &
 Point torques are zero. 
 &
 $\vx \times \vfext = \vzero$
 \\ 
\end{tabular}
\ee
\begin{Hresult}
The Cauchy stress tensor $\vsigma(\vx, t)$ is symmetric, thus
 $\vsigma = \vsigma\transpose$.
\end{Hresult}

\section{Conservation of Energy}
\be
\begin{tabular}{cl|l}
 &
 {\bf Assumption}
 &
 {\bf Expression} \\ \hline \hline
 (1) 
 &
 Energy is neither created nor destroyed.
 &
  
 \\ 
\end{tabular}
\ee

